\documentclass[12pt, oneside]{article}

% ── Packages ──────────────────────────────────────────────────────────────
\usepackage[T1]{fontenc}
\usepackage[utf8]{inputenc}
\usepackage[a4paper, margin=1in, top=1.1in, bottom=1.1in]{geometry}
\usepackage{times}
\usepackage{graphicx}
\usepackage{booktabs}
\usepackage{array}
\usepackage{tabularx}
\usepackage{longtable}
\usepackage{multirow}
\usepackage{microtype}
\usepackage{enumitem}
\usepackage{caption}
\usepackage{fancyhdr}
\usepackage{titlesec}
\usepackage{hyperref}
\usepackage{amsmath}
\usepackage{setspace}
\usepackage{parskip}
\usepackage{tocloft}
\usepackage{float}   % for [H] placement

% ── Spacing ───────────────────────────────────────────────────────────────
\setstretch{1.15}
\setlength{\parskip}{6pt}
\setlength{\parindent}{0pt}

% ── Hyperref ──────────────────────────────────────────────────────────────
\hypersetup{
  colorlinks=true,
  linkcolor=black,
  citecolor=black,
  urlcolor=blue,
  pdftitle={CSAS Technical Report --- Workshops 1 and 2},
  pdfauthor={Garzon, Quintero, Garrido, Cusba}
}

% ── Section formatting ────────────────────────────────────────────────────
\titleformat{\section}{\large\bfseries}{}{0em}{\thesection\quad}
\titleformat{\subsection}{\normalsize\bfseries}{}{0em}{\thesubsection\quad}
\titleformat{\subsubsection}{\normalsize\itshape}{}{0em}{\thesubsubsection\quad}
\titlespacing{\section}{0pt}{18pt}{6pt}
\titlespacing{\subsection}{0pt}{12pt}{4pt}
\titlespacing{\subsubsection}{0pt}{8pt}{2pt}

% ── Header / footer ───────────────────────────────────────────────────────
\pagestyle{fancy}
\fancyhf{}
\fancyhead[L]{\small\textit{CSAS --- Technical Report}}
\fancyhead[R]{\small\textit{Systems Analysis \& Design, 2026-I}}
\fancyfoot[C]{\small\thepage}
\renewcommand{\headrulewidth}{0.4pt}

% ── Caption style ─────────────────────────────────────────────────────────
\captionsetup{font=small, labelfont=bf, labelsep=period, skip=4pt}

% ── Column types ──────────────────────────────────────────────────────────
\newcolumntype{L}[1]{>{\raggedright\arraybackslash}p{#1}}
\newcolumntype{C}[1]{>{\centering\arraybackslash}p{#1}}

% ── No number on References ───────────────────────────────────────────────
\renewcommand{\refname}{References}

% =============================================================================
\begin{document}
% =============================================================================

% ── Title Page ───────────────────────────────────────────────────────────
\begin{titlepage}
\centering
\vspace*{1cm}

{\large\bfseries Universidad Distrital Francisco Jos\'e de Caldas}\\[4pt]
{\normalsize Computer Engineering Program --- School of Engineering}\\[2pt]
{\normalsize Systems Analysis \& Design --- Semester 2026-I}\\[2pt]
{\normalsize Eng. Carlos Andr\'es Sierra, M.Sc.}

\vspace{1.8cm}
\rule{\textwidth}{1pt}\\[12pt]
{\LARGE\bfseries Community Security Alert System (CSAS)\\[8pt]
Technical Report}\\[8pt]
{\large Workshops No.~1 \& No.~2 --- Systems Analysis and Design}\\[4pt]
\rule{\textwidth}{1pt}

\vspace{1.8cm}
{\normalsize
\begin{tabular}{ll}
\textbf{Felipe Jose Garzon Herrera}     & \textit{Code: 20251020132} \\[6pt]
\textbf{Juan Esteban Quintero Gordillo} & \textit{Code: 20251020137} \\[6pt]
\textbf{Henry Samuel Garrido Medina}   & \textit{Code: 20251020125} \\[6pt]
\textbf{Gabriel Mateo Cusba Marin}     & \textit{Code: 20251020128} \\
\end{tabular}}

\vspace{1.8cm}
{\normalsize Bogot\'a D.C., Colombia --- March 2026}

\vfill
{\small\textit{Submitted in partial fulfillment of the requirements for the\\
Systems Analysis \& Design course, Semester 2026-I.}}
\end{titlepage}

% ── Front matter ─────────────────────────────────────────────────────────
\pagenumbering{roman}

\section*{Abstract}
\addcontentsline{toc}{section}{Abstract}

Campus security at Universidad Distrital Francisco Jos\'e de Caldas is impaired by
three compounding problems identified through systematic primary data collection:
a persistent under-reporting culture attributable to the high friction of current
reporting channels, a spatial and temporal concentration of incidents during evening
and night hours in specific campus zones, and the absence of a centralized real-time
incident management tool that forces security staff to rely on phone calls and walk-up
reports, resulting in average response delays of approximately 11.4~minutes. This
report presents the Community Security Alert System (CSAS), a crowd-sourced
mobile and web platform that addresses these gaps through a microservices-based,
event-driven architecture integrating a hybrid AI-and-human verification pipeline,
geofenced alert dispatch with zone-weighted and time-weighted priority logic, and a
multi-channel notification engine. The proposed design targets a reduction of emergency
response time to under 5~minutes, an increase of incident recording completeness
from 41\% to above 85\%, and a community adoption rate of at least 60\%.

\vspace{6pt}
\noindent\textbf{Keywords:} campus security, crowd-sourced incident reporting,
microservices architecture, real-time notifications, geofenced alerts,
systems analysis, systems design, Universidad Distrital.

\newpage
\tableofcontents
\newpage
\listoffigures
\addcontentsline{toc}{section}{List of Figures}
\listoftables
\addcontentsline{toc}{section}{List of Tables}
\newpage

% ── Main matter ──────────────────────────────────────────────────────────
\pagenumbering{arabic}
\setcounter{page}{1}

% =============================================================================
\section{Introduction}
% =============================================================================

Campus safety is a persistent challenge for universities in urban Latin American
environments. At Universidad Distrital Francisco Jos\'e de Caldas, a public university
located in Bogot\'a D.C., the absence of a centralized digital reporting tool means that
security incidents are communicated through informal channels --- phone calls,
WhatsApp groups, and walk-up reports to security guards --- introducing delays, data
loss, and a systemic inability to detect spatial or temporal patterns. The average time
from incident occurrence to security team notification has been measured at approximately
11.4~minutes, a figure that represents a significant risk during fast-developing
emergencies such as assaults or medical events.

Several platforms have attempted to address similar challenges in the Bogot\'a
metropolitan area. The Polic\'ia Nacional's Vigila App integrates directly with law
enforcement but relies on the user placing a phone call, producing slow response cycles.
Neighborhood WhatsApp networks achieve high adoption rates due to their familiarity
but operate without any verification mechanism, systematically generating false alarms
and panic. SIURE~UD, the university's own internal incident reporting system,
provides an existing institutional process but operates exclusively through email and
phone, offering no mobile interface or real-time notification capability~\cite{sierra2026}.
These precedents establish a clear gap: no existing platform combines low-friction
mobile reporting, real-time geofenced notification, and verified alert quality in an
institutional campus context.

The Community Security Alert System (CSAS) is proposed to fill this gap. It is
designed as a crowd-sourced platform in which students, faculty, and staff act as
active sensors for community safety, submitting geolocated incident reports through a
mobile interface while a hybrid AI-and-human verification pipeline ensures signal
quality before alerts are broadcast to the affected campus community. This report
integrates two phases of the project: Workshop No.~1, which presents the empirical
systems analysis conducted to characterize the problem, and Workshop No.~2, which
presents the formal systems design derived from those findings.

% =============================================================================
\section{Literature Review}
% =============================================================================

The theoretical foundation of this project draws from established systems thinking and
engineering frameworks. Checkland's soft systems methodology~\cite{checkland1981}
provides the conceptual basis for treating campus security as a sociotechnical system
in which human behavior, institutional processes, and technology interact in complex,
non-linear ways. Meadows~\cite{meadows2008} offers the analytical vocabulary for
understanding the feedback dynamics that govern adoption and alert quality in a
crowd-sourced platform: the reinforcing loops that amplify participation and the
balancing loops that must be calibrated to prevent alert fatigue. The ISO/IEC/IEEE
15288:2015 standard~\cite{iso15288} provides the life-cycle process framework within
which the CSAS analysis and design activities are situated. Creswell and
Creswell~\cite{creswell2018} inform the mixed-methods data collection strategy
employed in Workshop No.~1, combining structured interviews, surveys, direct
observation, document analysis, benchmarking, and environmental monitoring into a
multi-method design that strengthens the validity of findings through methodological
triangulation. From the software architecture perspective, Bass, Clements, and
Kazman~\cite{bass2013} ground the design decisions related to modularity, scalability,
and separation of concerns that motivate the microservices decomposition. Dragoni et
al.~\cite{dragoni2017} provide a contemporary treatment of microservices as an
architectural style, validating its suitability for systems with heterogeneous load
profiles across subsystems --- precisely the case in the CSAS, where verification,
notification, and analytics exhibit independent scaling requirements.

% =============================================================================
\section{Objectives}
% =============================================================================

\subsection{General Objective}

To design a community-driven, real-time incident reporting and alert system for the
Universidad Distrital Francisco Jos\'e de Caldas campus that demonstrably reduces
emergency response times and improves incident recording completeness relative to the
current state.

\subsection{Specific Objectives}

\begin{enumerate}[leftmargin=*, itemsep=4pt]
\item To characterize the current campus security posture through primary data
collection, identifying key gaps, behavioral parameters, and system boundaries
(Workshop No.~1).
\item To translate the empirical findings of the systems analysis into a formal set of
functional and non-functional requirements explicitly traceable to observed problems.
\item To design a microservices-based architecture that addresses the identified gaps
through modular, independently scalable components.
\item To define a hybrid AI-and-human verification pipeline that balances alert latency
with alert signal quality.
\item To specify a technical implementation strategy --- including technology stack,
implementation phases, and key performance indicators --- as a foundation for future
development.
\end{enumerate}

% =============================================================================
\section{Scope}
% =============================================================================

The CSAS scope is defined by a boundary that separates the platform's direct
responsibilities from its operational environment. Table~\ref{tab:scope} formalizes
this boundary. The system covers the campus of Universidad Distrital Francisco
Jos\'e de Caldas and its immediate surrounding streets, and serves six stakeholder
groups: students (primary reporters and alert recipients), professors and administrative
staff (secondary reporters), campus security staff (verifiers and response coordinators),
local authorities (recipients of escalated verified incidents), system administrators
(platform maintainers), and university management (consumers of analytics and
security reports).

\begin{table}[H]
\caption{System Boundary Definition --- CSAS}
\label{tab:scope}
\renewcommand{\arraystretch}{1.2}
\centering
\begin{tabular}{L{6.5cm} L{6.5cm}}
\toprule
\textbf{Inside the system (scope)} & \textbf{Outside the system (environment)} \\
\midrule
Incident reporting and management platform & University physical infrastructure \\
Crowd-sourced verification logic & Local law enforcement (Police / EMS) \\
Real-time push notification engine & Public telecommunications and ISP networks \\
User and incident databases & Non-registered external visitors \\
Emergency coordination dashboard & Campus security physical hardware (CCTV, etc.) \\
Analytics and heatmap generation & University academic or administrative systems \\
\bottomrule
\end{tabular}
\end{table}

% =============================================================================
\section{Assumptions}
% =============================================================================

The analysis and design documented in this report rest on three critical assumptions
that have not yet been validated under real operating conditions and that represent the
primary behavioral and technical risks for the implementation phase.

\textbf{Willingness to report.} It is assumed that students and staff will submit
incident reports if the reporting process is made sufficiently simple and low-friction.
Workshop No.~1 survey data provides directional support: all participants who had
personally experienced a campus incident stated they would use the application.
However, this preference expressed under survey conditions has not been confirmed in
real operational settings. This assumption is designated for validation through user
acceptance testing in Phase~3.

\textbf{Verification latency tolerance.} The design assumes that the verification
mechanism can process incoming reports and produce a verified alert within a
45-second latency budget for standard-severity incidents. The validity of this
assumption depends entirely on implementation-specific decisions --- including AI
model performance, confirmation threshold configuration, and cloud infrastructure
capacity --- that fall outside the scope of this design document and must be validated
through load testing in Phase~4.

\textbf{Mobile connectivity coverage.} The system assumes reliable mobile network
and WiFi connectivity across all campus zones. The direct observation phase of
Workshop No.~1 did not include a systematic assessment of signal coverage. Dead
zones in low-connectivity areas could materially degrade alert delivery guarantees
and represent a prerequisite risk for infrastructure planning.

% =============================================================================
\section{Limitations}
% =============================================================================

Several limitations constrain the reach of the conclusions drawn in this report. These
are acknowledged not as weaknesses that invalidate the work, but as boundaries that
define the scope of the evidence base and establish a research agenda for subsequent
phases.

The student survey conducted in Workshop No.~1 ($n = 20$) is too small for statistical
generalization to the full university population. The 35\% incident prevalence estimate
is a meaningful directional indicator but cannot be treated as a representative rate
without a larger, stratified sample. The direct observation phase covered only four
campus zones across eight sessions; indoor corridors, cafeterias, and transit areas
between buildings were not observed, leaving potential blind spots in the risk map.
The document analysis was limited by restricted access to institutional security records,
which lacked the zone-specific and time-specific breakdowns needed for granular
historical validation. The weather correlation finding --- both suspicious events observed
during the study coincided with rainy nights --- is based on a one-week window and is
insufficient to establish causality or robust correlation. Finally, no implementation or
testing has been conducted at this stage; all performance and reliability figures presented
in the design are targets, not measured outcomes.

Despite these limitations, the data collected is considered sufficient to guide the design
phase. The most important findings --- night-time incident concentration, low reporting
levels, and preference for a mobile channel --- appear independently across interviews,
surveys, and direct observation. When three independent methods converge on the same
conclusion, confidence in that conclusion is significantly higher than any single method
could provide.

% =============================================================================
\section{Methodology}
% =============================================================================

\subsection{Systems Analysis Approach (Workshop No.~1)}

The systems analysis phase employed a mixed-methods approach~\cite{creswell2018}
that combined six complementary data collection methods to characterize the campus
security problem from multiple independent perspectives. Table~\ref{tab:dcplan}
summarizes the six methods applied.

\begin{table}[H]
\caption{Data Collection Plan --- Six Methods Applied to the CSAS}
\label{tab:dcplan}
\renewcommand{\arraystretch}{1.2}
\centering
\begin{tabular}{L{2.2cm} L{2.0cm} L{2.4cm} L{1.4cm} L{4.5cm}}
\toprule
\textbf{Method} & \textbf{Tool} & \textbf{Sample} & \textbf{Dur.} & \textbf{Objective} \\
\midrule
Interviews & Structured guide & 5 security, 2 admin staff & 1 week & Understand current workflows and identify unmet needs. \\
\addlinespace
Surveys & Google Forms & 20 students & 2 weeks & Assess safety perception, incident experience, and app adoption intent. \\
\addlinespace
Direct observation & Record form & 4 campus zones & 1 week & Record activity level and suspicious events by time and zone. \\
\addlinespace
Document analysis & Institutional reports & Campus records & 3 days & Identify historical incident patterns and cross-validate field findings. \\
\addlinespace
Benchmarking & Comparative analysis & 3 apps in Bogot\'a & 1 week & Extract design insights from existing platforms. \\
\addlinespace
Weather monitoring & Weather apps & Teusaquillo locality & 1 week & Explore correlation between adverse weather and incident occurrence. \\
\bottomrule
\end{tabular}
\end{table}

\textbf{Structured interviews} were conducted with five security staff members and two
administrative staff members. Three main problems emerged consistently: communication
times are slow (the response team may take several minutes to be notified); there is no
centralized reporting tool; and staff need faster notifications that include the exact
location of incidents, not just verbal descriptions.

\textbf{Student survey} ($n = 20$) measured safety perception, incident experience,
incident type, willingness to use a reporting application, and time of day. Seven of the
20 students (35\%) reported having experienced or witnessed some type of incident on
or around campus. All reported incidents occurred during afternoon and night hours.
100\% of students who experienced an incident stated they would use the alert
application, compared to 69\% among those who had not.

\textbf{Direct observation} was conducted across four campus zones over eight sessions
at different times of day. Of the eight sessions, only two recorded suspicious events,
and both occurred in the parking lot at night. The parking lot at night was identified
as the highest-risk observed zone, with poor lighting and low foot traffic as contributing
factors. The suspicious events observed were not reported by any other witness, directly
reinforcing the under-reporting hypothesis.

\textbf{Document analysis} of institutional security records confirmed that formal
records show far fewer incidents than actually occur according to survey and observation
data. The most documented incident types are theft and unauthorized access.

\textbf{Benchmarking} evaluated three existing security alert platforms in Bogot\'a.
Results are presented in Appendix~\ref{app:bench}.

\textbf{Weather monitoring} tracked conditions in the Teusaquillo locality for one
week. Both suspicious events coincided with rainy nights, motivating a weather API
integration in the system design.

Figure~\ref{fig:context} shows the stakeholder relationships and information flows.
Figure~\ref{fig:nodegraph} shows the full system architecture as a node relationship
graph from Workshop No.~1.

\begin{figure}[H]
\centering
\includegraphics[width=0.82\textwidth]{fig_context.png}
\caption{Stakeholder relationships and information flows in the Community Security
Alert System (Workshop No.~1).}
\label{fig:context}
\end{figure}

\begin{figure}[H]
\centering
\includegraphics[width=0.82\textwidth]{fig_nodegraph.png}
\caption{Node relationship graph of the Community Security Alert System. Node size
reflects relative centrality; the REST API is the highest-connectivity node
(Workshop No.~1).}
\label{fig:nodegraph}
\end{figure}

\subsection{System Architecture Design (Workshop No.~2)}

The systems design phase translated the empirical findings of Workshop No.~1 into a
formal architecture using an iterative, requirements-driven approach consistent with
the ISO/IEC/IEEE 15288:2015 systems life cycle~\cite{iso15288}. Functional and
non-functional requirements were derived directly from identified gaps, with explicit
traceability between each requirement and the Workshop No.~1 finding that motivated
it. The architectural style was selected based on the non-uniform load profiles of the
CSAS subsystems and the need for independent scalability, motivating a
microservices-based, event-driven decomposition~\cite{dragoni2017, bass2013}.

The architecture is organized across five layers --- Presentation, Application, Service,
Data, and Integration. The five core microservices (Incident, Verification, Dispatcher,
User, and Analytics) communicate asynchronously through a message broker, ensuring
loose coupling between the ingestion and processing stages. Figure~\ref{fig:arch}
shows the architecture diagram. Figure~\ref{fig:flow} shows the main process flow
from incident submission to alert dispatch.

\begin{figure}[H]
\centering
\includegraphics[width=0.82\textwidth]{fig1_architecture.png}
\caption{CSAS system architecture organized across five layers. The Service Layer
hosts the five core microservices; external integrations (Authorities, SIURE~UD,
SMS Gateway) connect through the Integration Layer (Workshop No.~2).}
\label{fig:arch}
\end{figure}

\begin{figure}[H]
\centering
\includegraphics[width=0.55\textwidth]{fig2_processflow.png}
\caption{Main process flow diagram. An incident submitted by a user passes through
the Verification Service; valid incidents trigger alert dispatch while invalid reports
are rejected and logged (Workshop No.~2).}
\label{fig:flow}
\end{figure}

\subsection{Technology Stack}

Table~\ref{tab:stack} summarizes the recommended technology stack. Each choice is
justified with reference to a specific functional or non-functional requirement.

\begin{table}[H]
\caption{Recommended Technology Stack and Justification}
\label{tab:stack}
\renewcommand{\arraystretch}{1.2}
\centering
\begin{tabular}{L{3.0cm} L{3.0cm} L{7.0cm}}
\toprule
\textbf{Component} & \textbf{Technology} & \textbf{Justification} \\
\midrule
Mobile application  & React Native       & Single codebase for Android/iOS; supports $\leq$3-interaction reporting (NFR-06). \\
API gateway         & Node.js / Nginx    & High I/O throughput; handles authentication, rate limiting, and routing. \\
Backend microservices & Node.js          & Asynchronous, event-driven runtime supports the latency target of NFR-01 ($<$30\,s). \\
Verification Engine & Python / FastAPI   & Python's ML ecosystem integrates the AI plausibility scorer; FastAPI provides high-performance async endpoints. \\
Message broker      & RabbitMQ           & Priority queue support; decouples ingestion from verification to mitigate volume spikes. \\
Primary database    & PostgreSQL + PostGIS & Native geospatial queries for zone-weighted routing and heatmap generation (FR-06, FR-09). \\
Push notifications  & Firebase FCM       & Cross-platform delivery; supports the 99.9\% delivery target of NFR-02. \\
SMS fallback        & Twilio API         & Guaranteed delivery for users without smartphone access; supports equity of access. \\
Infrastructure      & AWS / Azure (K8s)  & Container orchestration enables per-service auto-scaling (NFR-03, NFR-05). \\
\bottomrule
\end{tabular}
\end{table}

% =============================================================================
\section{Requirements}
% =============================================================================

\subsection{Functional Requirements}

Table~\ref{tab:fr} presents the ten functional requirements derived directly from
Workshop No.~1. Each requirement is explicitly traceable to at least one empirical finding.

\begin{table}[H]
\caption{Functional Requirements and Workshop No.~1 Traceability}
\label{tab:fr}
\renewcommand{\arraystretch}{1.2}
\centering
\begin{tabular}{L{1.3cm} L{3.2cm} L{8.5cm}}
\toprule
\textbf{ID} & \textbf{Requirement} & \textbf{W1 Rationale} \\
\midrule
FR-01 & Incident reporting via GPS, type, and description & High reporting friction is the primary cause of under-reporting observed across all data collection methods. \\
FR-02 & Geofenced alerts within 500\,m of verified incidents & Spatial concentration of incidents in specific zones requires targeted, not broadcast, notifications. \\
FR-03 & AI plausibility scorer for initial validation & Data fragmentation in current manual records increases false positives; automated filtering is needed. \\
FR-04 & Administrative dashboard for security staff & Fragmented phone and walk-up reporting creates the 11.4-minute delay; a centralized view is required. \\
FR-05 & Crowd-sourced confirmation and flagging & 100\% of incident-experienced survey participants would use the platform, indicating community validation potential. \\
FR-06 & Automatic GPS capture on report submission & Spatial precision requirements identified during four-zone observation sessions. \\
FR-07 & Persistent incident log for pattern analysis & Currently only 41\% of events are recorded; historical continuity is absent. \\
FR-08 & One-touch emergency SOS button & Targets a more than 50\% reduction from the 11.4-minute baseline response time. \\
FR-09 & Periodic security heatmaps for management & Benchmarking revealed absence of data-driven oversight in comparable platforms. \\
FR-10 & Image and video attachments to reports & Improves evidence quality and verification reliability for ambiguous reports. \\
\bottomrule
\end{tabular}
\end{table}

\subsection{Non-Functional Requirements}

Table~\ref{tab:nfr} defines the quality constraints that bound the system's behavior
independently of specific functions. Each NFR establishes a measurable threshold against
which the implemented system will be evaluated.

\begin{table}[H]
\caption{Non-Functional Requirements and Target Metrics}
\label{tab:nfr}
\renewcommand{\arraystretch}{1.2}
\centering
\begin{tabular}{L{1.3cm} L{2.8cm} L{9.0cm}}
\toprule
\textbf{ID} & \textbf{Category} & \textbf{Requirement} \\
\midrule
NFR-01 & Latency       & End-to-end latency from submission to broadcast must be $<$30\,s for 95\% of incidents. \\
NFR-02 & Reliability   & Notification delivery success $>$99.9\% across push and SMS channels. \\
NFR-03 & Availability  & System uptime $\geq$99.9\%; guaranteed during high-risk hours (18:00--22:00). \\
NFR-04 & Privacy/Security & All transmissions via TLS~1.3; full compliance with Ley~1581 de~2012~\cite{ley1581}. \\
NFR-05 & Scalability   & API gateway must sustain a 300\% concurrent request surge without performance degradation. \\
NFR-06 & Usability     & Incident reporting workflow completable in $\leq$3 interactions. \\
NFR-07 & Maintainability & Any microservice must be updatable or replaceable independently. \\
NFR-08 & Data Integrity & No alert shall broadcast without $\geq$2 community confirmations or 1 administrative approval. \\
\bottomrule
\end{tabular}
\end{table}

% =============================================================================
\section{Implementation Plan}
% =============================================================================

As this report covers the analysis and design phases only, no implemented system
exists at this stage. This section presents the planned implementation strategy,
the key performance indicators against which the future system will be evaluated,
and the testing approach that will validate the design decisions.

\subsection{Implementation Phases}

The project follows four sequential phases aligned with the system life cycle.

\textbf{Phase~1 --- Planning and Architecture (Month~1).} Finalizes the system
architecture, validates all requirements against stakeholder needs, and establishes the
development environment. \textit{Success criterion:} all FR and NFR are documented,
agreed upon, and traceable to Workshop No.~1 findings. \textit{Deliverable:}
architecture diagrams and validated requirements specification.

\textbf{Phase~2 --- Development (Month~2).} Implements the frontend (React Native),
backend microservices (Node.js, Python), database schema (PostgreSQL~+~PostGIS), and
messaging infrastructure (RabbitMQ). \textit{Success criterion:} an end-to-end prototype
in which a user can submit a geolocated incident report (FR-01) and receive a geofenced
alert (FR-02). \textit{Deliverable:} functional prototype.

\textbf{Phase~3 --- Testing and Validation (Month~3).} Executes the full testing
strategy (see Section~\ref{sec:testing}). \textit{Success criterion:} system meets
NFR-01 ($<$30\,s latency) and NFR-02 (99.9\% delivery) under simulated load.
\textit{Deliverable:} test reports with pass/fail evidence for all defined test cases.

\textbf{Phase~4 --- Deployment and Monitoring (Month~4).} Deploys to the production
cloud environment, activates external integrations (SIURE~UD, SMS gateway), and
establishes live monitoring dashboards. \textit{Success criterion:} 99.9\% uptime
under real-world conditions. \textit{Deliverable:} deployed system with operational
KPI dashboard.

\subsection{Key Performance Indicators}

Table~\ref{tab:kpis} defines the measurable targets against which the implemented
system will be evaluated. Each KPI is aligned with a specific requirement from
Section~\ref{sec:req}.

\begin{table}[H]
\caption{Key Performance Indicators, Targets, and Measurement Methods}
\label{tab:kpis}
\renewcommand{\arraystretch}{1.2}
\centering
\begin{tabular}{L{4.2cm} L{3.0cm} L{3.4cm} L{1.8cm}}
\toprule
\textbf{KPI} & \textbf{Target} & \textbf{Measurement} & \textbf{Freq.} \\
\midrule
Alert latency (end-to-end) & $<$30\,s (NFR-01)  & System event logs        & Daily      \\
Notification delivery rate & $>$99.9\% (NFR-02) & FCM / Twilio receipts    & Daily      \\
System availability        & 99.9\% (NFR-03)    & Uptime monitoring        & Continuous \\
Incident recording rate    & $>$85\% of events  & Cross-ref.\ staff records & Monthly   \\
Community adoption rate    & $>$60\% of students & Monthly active users     & Monthly    \\
Verification latency       & $<$60\,s           & Event broker timestamps  & Daily      \\
Alert accuracy rate        & $\leq$5\% false alarms & Post-incident review  & Weekly     \\
Report submission time     & $\leq$2 min (NFR-06) & UAT session logs       & Per UAT    \\
\bottomrule
\end{tabular}
\end{table}

\subsection{Testing Strategy}
\label{sec:testing}

\textbf{Unit testing} will validate individual components within each microservice
(authentication module, incident submission endpoint, notification dispatch function,
AI scorer). Each service shall maintain at least 80\% code coverage.

\textbf{Integration testing} will verify the complete event chain --- Incident Service
$\rightarrow$ Verification Engine $\rightarrow$ Alert Dispatcher --- using a staging
environment seeded with campus zone data and simulated user accounts.

\textbf{Load testing} will evaluate system behavior under a 300\% traffic surge
consistent with NFR-05. The acceptance criterion is zero timeout errors and
end-to-end latency below 30\,s at the 95th percentile.

\textbf{User acceptance testing (UAT)} will engage a pilot group of 50 students and
5 security staff members. The primary acceptance criterion is that at least 85\% of
participants complete their first incident report submission within two minutes.

\textbf{Security testing} will verify TLS~1.3 encryption, authentication mechanisms,
access control policies, and compliance with Ley~1581 de~2012~\cite{ley1581},
addressing NFR-04.

% =============================================================================
\section{Discussion}
% =============================================================================

\subsection{Operational Gap Analysis}

The primary motivation for the CSAS is a measurable gap between the current state
of campus security and the operational targets the platform is designed to achieve.
Table~\ref{tab:gaps} quantifies the most critical dimensions of this gap as established
in Workshop No.~1.

\begin{table}[H]
\caption{Operational Gap Analysis: Current State vs.\ Design Targets}
\label{tab:gaps}
\renewcommand{\arraystretch}{1.2}
\centering
\begin{tabular}{L{4.5cm} L{3.0cm} L{2.5cm} L{2.5cm}}
\toprule
\textbf{Dimension} & \textbf{Current State} & \textbf{Target} & \textbf{Gap} \\
\midrule
Emergency response time      & $\sim$11.4\,min   & $<$5\,min    & $-$6.4\,min \\
Incident recording completeness & $\sim$41\%     & $>$85\%      & $+$44\,pp   \\
Community app adoption       & $\sim$12\%        & $\geq$60\%   & $+$48\,pp   \\
Alert delivery success rate  & $\sim$87\% (SMS/call) & $>$99.9\% & $+$12.9\,pp \\
\bottomrule
\end{tabular}
\end{table}

These gaps are not exclusively technical. Closing the adoption gap from 12\% to 60\%,
for instance, requires not only a functional mobile application but a sustained community
engagement strategy and explicit institutional endorsement from university management
--- factors that lie at the boundary of the system scope but are critical to real-world
effectiveness.

\subsection{Risk and Complexity Management}

Workshop No.~1 identified five operational parameters with the strongest influence on
CSAS behavior. The design incorporates a specific mitigation mechanism for each.
A sharp increase in report volume is mitigated by the asynchronous message broker
that decouples ingestion from verification, preventing the Verification Engine from
becoming a bottleneck. Low user participation is addressed by the low-friction mobile
interface (FR-01, $\leq$3 interactions) and the community engagement module planned
for Phase~3. Night-time incident concentration is handled by time-weighted priority
logic that automatically elevates the urgency of off-hours alerts and activates an
on-call escalation path for security staff. High-risk zone clustering --- particularly
the parking lot, identified as the highest-risk zone across all observation sessions ---
is managed by zone-based rate limiting and spatial clustering of co-located simultaneous
reports. Adverse weather correlation, while preliminary (based on only two co-occurring
observations), will be addressed through real-time weather API integration that adjusts
monitoring sensitivity in high-risk zones when conditions deteriorate.

\subsection{Feedback Dynamics}

The CSAS exhibits two interacting feedback dynamics~\cite{meadows2008}. The
\textit{reinforcing loop} operates as follows: useful alerts increase user vigilance,
which increases reporting, which generates more alerts and further improves system
effectiveness as adoption grows. The community engagement module targets the critical
adoption threshold of approximately 20--25\% of the campus population, the point at
which report density becomes sufficient to sustain the verification mechanism without
heavy dependence on security staff moderation. Below this threshold, the system relies
more strongly on staff-initiated verification. The \textit{balancing loop} operates in
the opposite direction: excessive alert volume triggers alert fatigue, reducing user
engagement and system effectiveness. Two mechanisms work in concert to prevent this
degradation: the alert fatigue dampening module in the User Service, which progressively
reduces non-critical notification frequency for users with high dismissal rates, and the
Verification Engine's quality filters, which prevent low-confidence reports from reaching
the dispatcher. Calibrating the confirmation threshold --- fast enough for genuine
emergencies, strict enough to suppress noise --- is the most critical design decision
for the implementation phase.

\subsection{Ethical Considerations}

The CSAS collects three categories of personal data --- registration information, real-time
location data, and incident media --- each subject to distinct obligations under Ley~1581
de~2012, Colombia's general personal data protection regime~\cite{ley1581}. The design
incorporates four binding obligations: explicit informed consent at registration specific
to the safety purpose; purpose limitation (data may not be repurposed or transferred to
third parties); user rights to access, correct, and delete personal data; and a defined
retention policy (maximum 12~months for incident records before deletion or irreversible
anonymization). For sensitive incident types such as harassment or threats, an anonymous
reporting option encrypts the reporter's identity at rest, accessible only to system
administrators under a documented access-controlled protocol. A formal data processing
agreement with the university's legal counsel is a prerequisite to production deployment.
Equity of access is addressed through the multi-channel notification architecture: the
SMS gateway (NFR-02) ensures that users without smartphones or mobile data access can
both receive alerts and submit reports, which is particularly relevant given socioeconomic
diversity among the student population.

% =============================================================================
\section{Conclusions}
% =============================================================================

This report has documented the complete analysis and design phases of the Community
Security Alert System for Universidad Distrital Francisco Jos\'e de Caldas. The systems
analysis conducted in Workshop No.~1, through six independent data collection methods
whose findings consistently converged on the same three core problems, established an
empirical foundation that is robust for a project at this stage. The convergence of
interviews, student surveys, direct observation, document analysis, benchmarking, and
weather monitoring on the same conclusions --- under-reporting, night-time and spatial
incident concentration, and the absence of a centralized tool --- provides high confidence
that the problem characterization is accurate and that design decisions grounded in these
findings are well-motivated.

The architecture proposed in Workshop No.~2 directly addresses each identified gap.
The low-friction mobile interface and community engagement module target the
under-reporting culture. The adaptive zone-weighted and time-weighted dispatch logic
of the Alert Dispatcher addresses the non-uniform distribution of incidents. The
operational performance gaps are formalized as measurable KPI targets in
Table~\ref{tab:kpis} that will enable objective evaluation of the implemented system.
Every major design decision is traceable to a specific empirical finding from
Workshop No.~1. Five systems engineering principles guided the design throughout:
modularity (independent microservices with well-defined interfaces), scalability
(per-component auto-scaling), maintainability (containerized CI/CD deployment),
reliability (redundancy mechanisms and offline queuing), and privacy
(compliance-by-design with Ley~1581 de~2012).

Three open assumptions require validation in the implementation phase: user willingness
to report under real operating conditions, verification latency within the 45-second
budget, and mobile connectivity coverage across all campus zones. Future work should
expand the student survey to a statistically representative sample (minimum $n = 200$,
stratified by faculty and time of campus use), conduct a systematic campus connectivity
audit before finalizing the notification delivery architecture, and establish formal
data-sharing and data processing agreements with SIURE~UD and university legal counsel
before production deployment.

% ── References ───────────────────────────────────────────────────────────
\newpage
\begin{thebibliography}{9}

\bibitem{sierra2026}
C.~A.~Sierra,
\textit{TeamWork as Computer Engineers --- CS Collaboration Guidelines}.
Bogot\'a: Universidad Distrital Francisco Jos\'e de Caldas, 2026.

\bibitem{checkland1981}
P.~Checkland,
\textit{Systems Thinking, Systems Practice}.
New York: John Wiley \& Sons, 1981.

\bibitem{meadows2008}
D.~H.~Meadows,
\textit{Thinking in Systems: A Primer}.
White River Junction, VT: Chelsea Green Publishing, 2008.

\bibitem{iso15288}
ISO/IEC/IEEE,
\textit{Systems and Software Engineering --- System Life Cycle Processes},
ISO/IEC/IEEE Std.~15288:2015, 2015.

\bibitem{creswell2018}
J.~W.~Creswell and J.~D.~Creswell,
\textit{Research Design: Qualitative, Quantitative, and Mixed Methods Approaches},
5th~ed. Thousand Oaks, CA: SAGE Publications, 2018.

\bibitem{ley1581}
Congreso de Colombia,
\textit{Ley 1581 de 2012 --- R\'egimen General de Protecci\'on de Datos Personales},
2012.

\bibitem{bass2013}
L.~Bass, P.~Clements, and R.~Kazman,
\textit{Software Architecture in Practice}, 3rd~ed.
Boston: Addison-Wesley, 2013.

\bibitem{dragoni2017}
N.~Dragoni \textit{et al.},
``Microservices: Yesterday, Today, and Tomorrow,''
\textit{Present and Ulterior Software Engineering},
pp.~195--216, 2017.

\end{thebibliography}

% ── Appendices ───────────────────────────────────────────────────────────
\newpage
\appendix

\section{Raw Survey Data (Workshop No.~1)}
\label{app:survey}

Table~\ref{tab:survey} presents the complete raw data from the student survey ($n = 20$)
conducted during Workshop No.~1. Variables collected were: campus safety perception,
incident experience, incident type, willingness to use the alert application, and time of
day when incidents occurred.

\begin{table}[H]
\caption{Raw survey data ($n=20$). Variables: safety perception, incident experience,
incident type, app adoption intention, and time of day.}
\label{tab:survey}
\renewcommand{\arraystretch}{1.1}
\centering
\small
\begin{tabular}{C{0.7cm} C{1.7cm} C{2.0cm} C{2.8cm} C{2.2cm} C{2.0cm}}
\toprule
\textbf{ID} & \textbf{Feels Safe} & \textbf{Experienced} & \textbf{Incident Type} & \textbf{Would Use App} & \textbf{Time} \\
\midrule
1  & No  & Yes & Theft               & Yes   & Night     \\
2  & Yes & No  & None                & Yes   & Afternoon \\
3  & No  & Yes & Harassment          & Yes   & Night     \\
4  & Yes & No  & None                & Maybe & Morning   \\
5  & Yes & No  & None                & Maybe & Morning   \\
6  & Yes & No  & None                & Yes   & Afternoon \\
7  & No  & Yes & Theft               & Yes   & Night     \\
8  & Yes & No  & None                & Maybe & Morning   \\
9  & Yes & No  & None                & Maybe & Morning   \\
10 & Yes & No  & None                & Yes   & Afternoon \\
11 & No  & Yes & Suspicious Activity & Yes   & Night     \\
12 & Yes & No  & None                & Maybe & Morning   \\
13 & No  & Yes & Theft               & Yes   & Evening   \\
14 & Yes & No  & None                & Yes   & Afternoon \\
15 & No  & Yes & Harassment          & Yes   & Night     \\
16 & Yes & No  & None                & Maybe & Morning   \\
17 & Yes & No  & None                & Maybe & Night     \\
18 & Yes & No  & None                & Yes   & Afternoon \\
19 & No  & Yes & Theft               & Yes   & Night     \\
20 & Yes & No  & None                & Maybe & Morning   \\
\midrule
\multicolumn{2}{l}{\textbf{Experienced incident:}} &
\multicolumn{4}{l}{Yes: 7 (35\%) \quad No: 13 (65\%)} \\
\bottomrule
\end{tabular}
\end{table}

\section{Direct Observation Record (Workshop No.~1)}
\label{app:obs}

Table~\ref{tab:obs} presents the complete direct observation record from the eight
sessions conducted across four campus zones during Workshop No.~1.

\begin{table}[H]
\caption{Direct observation record --- 8 sessions across 4 campus zones. Variables:
activity level and presence of suspicious event.}
\label{tab:obs}
\renewcommand{\arraystretch}{1.15}
\centering
\begin{tabular}{C{1.8cm} L{3.2cm} C{2.2cm} C{2.8cm} C{2.8cm}}
\toprule
\textbf{Session} & \textbf{Location} & \textbf{Time} & \textbf{Activity Level} & \textbf{Suspicious Event} \\
\midrule
1 & Main Entrance  & Morning   & High   & No           \\
2 & Library Area   & Afternoon & High   & No           \\
3 & Parking Lot    & Night     & Medium & \textbf{Yes} \\
4 & Dormitory Area & Evening   & Medium & No           \\
5 & Main Entrance  & Afternoon & High   & No           \\
6 & Parking Lot    & Night     & Low    & \textbf{Yes} \\
7 & Library Area   & Morning   & Medium & No           \\
8 & Dormitory Area & Evening   & Medium & No           \\
\bottomrule
\end{tabular}
\end{table}

\section{Benchmarking of Existing Security Platforms}
\label{app:bench}

Table~\ref{tab:bench} presents the benchmarking comparison of three existing security
alert platforms operating in Bogot\'a conducted during Workshop No.~1.

\begin{table}[H]
\caption{Benchmarking of existing security alert platforms in Bogot\'a.}
\label{tab:bench}
\renewcommand{\arraystretch}{1.2}
\centering
\begin{tabular}{L{2.4cm} L{2.8cm} L{2.8cm} L{4.5cm}}
\toprule
\textbf{Platform} & \textbf{Strengths} & \textbf{Weaknesses} & \textbf{Insight for CSAS} \\
\midrule
Polic\'ia Nacional / Vigila App
& Direct police integration; national coverage
& Slow response; requires a phone call
& Authority integration model is valuable but must be faster and mobile-first. \\
\addlinespace
Neighborhood networks (WhatsApp)
& High adoption; instant notifications; familiar channel
& No verification; no records; generates panic
& Immediacy is a key design goal, but verification is non-negotiable. \\
\addlinespace
SIURE UD (internal)
& Exists; staff familiar; institutional endorsement
& Email and phone only; no mobile or real-time capability
& CSAS complements rather than replaces SIURE UD, building on institutional familiarity. \\
\bottomrule
\end{tabular}
\end{table}

\end{document}
